% !TeX program = xelatex
\documentclass[11pt,a4paper]{article}

% Компилировать с XeLaTeX или LuaLaTeX.
\usepackage[a4paper,margin=0.62in]{geometry}
\usepackage{fontspec}
\IfFontExistsTF{Lato}{\setmainfont{Lato}}{\setmainfont{Noto Sans}}
\usepackage{polyglossia}
\setdefaultlanguage{russian}
\usepackage{microtype}
\usepackage{tabularx}
\usepackage{array}
\usepackage{enumitem}
\usepackage{xcolor}
\usepackage{hyperref}
\usepackage{fancyhdr}
\usepackage{lastpage}

\definecolor{linkblue}{RGB}{20,33,110}
\hypersetup{
  colorlinks=true,
  urlcolor=linkblue,
  linkcolor=linkblue,
  citecolor=linkblue
}

\pagestyle{fancy}
\fancyhf{}
\renewcommand{\headrulewidth}{0pt}
\renewcommand{\footrulewidth}{0pt}
\cfoot{\small \thepage\ из \pageref{LastPage}}

\setlength{\parindent}{0pt}
\setlength{\parskip}{0pt}
\setlength{\tabcolsep}{0pt}
\setlist[itemize]{leftmargin=1.15em,itemsep=0.08em,topsep=0.15em,parsep=0pt,partopsep=0pt}

\newcolumntype{L}{>{\raggedright\arraybackslash}p{0.145\textwidth}}
\newcolumntype{R}{>{\raggedright\arraybackslash}p{0.835\textwidth}}

\newcommand{\cvrow}[2]{%
  \noindent
  \begin{minipage}[t]{0.145\textwidth}
    \vspace{0pt}%
    \raggedright
    \fontsize{8.5}{10}\selectfont
    \bfseries\MakeUppercase{#1}
  \end{minipage}%
  \hfill
  \begin{minipage}[t]{0.835\textwidth}
    \vspace{0pt}%
    #2
  \end{minipage}\par
}

\newcommand{\entry}[2]{%
  \begin{tabularx}{\linewidth}{@{}X >{\raggedleft\arraybackslash}p{0.28\linewidth}@{}}
    #1 & #2
  \end{tabularx}
}

\newcommand{\pub}[2]{%
  #1 #2\par
  \vspace{0.28em}
}

\newcommand{\eduentry}[2]{%
  \noindent
  \begin{tabularx}{\linewidth}
    {@{}X@{\hspace{0.8em}}
    >{\raggedleft\arraybackslash}p{0.30\linewidth}@{}}
    #1 & \small #2
  \end{tabularx}\par
  \vspace{0.15em}
}

\newcommand{\edudetail}[1]{%
  \noindent\hspace*{1.15em}$\diamond$ #1\par
}

\newcommand{\skillcell}[2]{%
  \textbf{#1}\par
  #2
}

\newcommand{\workentry}[2]{%
  \noindent
  \begin{tabularx}{\linewidth}
    {@{}X@{\hspace{0.8em}}
    >{\raggedleft\arraybackslash}p{0.30\linewidth}@{}}
    #1 & \small #2
  \end{tabularx}\par
  \vspace{0.15em}
}

\newcommand{\workproject}[1]{%
  \noindent
  \hangindent=2.30em
  \hangafter=1
  \hspace*{1.15em}\makebox[1.15em][l]{$\diamond$}#1\par
  \vspace{0.12em}
}

\newcommand{\workprojectdetail}[1]{%
  \noindent
  \hangindent=3.45em
  \hangafter=1
  \hspace*{2.30em}\makebox[1.15em][l]{--}#1\par
  \vspace{0.08em}
}

\begin{document}

% ── Шапка ─────────────────────────────────────────────
{\fontsize{24}{28}\selectfont\bfseries Артур Ясеновец}\par
\vspace{0.08em}
{\large\bfseries Инженер по информационной безопасности и прикладной криптографии}\par
\vspace{0.28em}
\href{mailto:archee.busy@gmail.com}{archee.busy@gmail.com}
\;\textbullet\;
\href{https://github.com/iasenovets}{GitHub}
\;\textbullet\;
\href{https://iasenovets.github.io/}{Портфолио}
\;\textbullet\;
\href{https://orcid.org/0009-0008-0008-3746}{ORCID}
\;\textbullet\;
Чунцин, Китай
\par
\vspace{0.22em}
\hrule
\vspace{0.55em}

% ── Профиль ────────────────────────────────────────────
\cvrow{Профиль}{%
  Инженер и исследователь в области информационной безопасности
  с фокусом на прикладную криптографию и технологии повышения
  конфиденциальности. Разрабатываю на \textbf{Go},
  \textbf{Python} и \textbf{TypeScript}; имею базовый опыт работы
  с \textbf{C++} и \textbf{Rust}. Занимаюсь формальным анализом
  безопасности криптографических протоколов и разработкой систем
  приватного поиска, извлечения и обработки данных на основе
  \textbf{гомоморфного шифрования}.
}

\vspace{0.50em}

\cvrow{Навыки}{%
  \begin{tabularx}{\linewidth}
    {@{}>{\raggedright\arraybackslash}X
    @{\hspace{1.2em}}
    >{\raggedright\arraybackslash}X@{}}

    \skillcell{Стандарты и регуляторика:}{%
      ГОСТ Р 34.10/34.11, ГОСТ 34.12/34.13;
      149-ФЗ, 152-ФЗ, 187-ФЗ, 63-ФЗ; FIPS 203--205}
    &
    \skillcell{Библиотеки и фреймворки:}{%
      Lattigo, Microsoft SEAL, OpenFHE, TenSEAL,
      feanor, fheanor}
    \\[0.20em]

    \skillcell{Технологии конфиденциальности:}{%
      Private Information Retrieval (PIR/CPIR),
      Searchable Symmetric Encryption (SSE)}
    &
    \skillcell{Системы и инфраструктура:}{%
      Linux, Docker, Git, Airflow, S3/MinIO,
      Hyperledger Fabric}
    \\[0.20em]

    \skillcell{Гомоморфное шифрование:}{%
      Brakerski--Gentry--Vaikuntanathan (BGV),
      Brakerski--Fan--Vercauteren (BFV),
      Cheon--Kim--Kim--Song (CKKS);
      анализ безоп-ти IND-CPA, IND-CCA2, EUF-CMA}
    &
    \skillcell{Инженерия безопасности:}{%
      YARA, Chainsaw, анализ вредоносного ПО и журналов,
      извлечение IoC, сопоставление TTP с MITRE ATT\&CK}
    \\[0.20em]

    \skillcell{Постквантовая криптография:}{%
      ML-KEM, ML-DSA, SLH-DSA;
      LWE/RLWE, Module-LWE, SIS/Module-SIS}
    &
    \skillcell{Языки:}{%
      русский --- родной;
      \href{https://iasenovets.github.io/files/toefl_8cd5f4a2b8ec8f34da49ea0faeef43f1.pdf}{английский --- C1};
    \href{https://admin.chinesetest.cn/queryScore.do?sid=c2tActR9Wb\%2704b861aefbdc5e9fd0b44afa717d08cbceb0720e582ce25f\%2724d7sWsQ5nVUQSwD}
        {китайский --- HSK 3}}
  \end{tabularx}
}

\vspace{0.50em}

\cvrow{Опыт работы}{%
  \workentry
    {\textbf{Магистрант-исследователь} @ Группа прикладной криптографии, CQUPT}
    {сент. 2024--наст. время}

  \workproject{\textbf{Приватное извлечение информации для Hyperledger Fabric}}
  \workprojectdetail{Спроектировал и реализовал на Go вычислительный PIR-протокол
  на основе BGV с использованием Lattigo, интегрировав приватное чтение из мирового
  состояния непосредственно в чейнкод Hyperledger Fabric.}
  \workprojectdetail{Разработал модель угроз с честным, но любопытным противником
  (HBC) и проанализировал конфиденциальность запросов на основе IND-CPA-безопасности
  BGV, включая безопасность повторных запросов и наблюдаемую утечку информации.}
  % \workprojectdetail{Построил экспериментальный стенд и оценил корректность,
  % задержку, объём коммуникации, требования к хранению и накладные расходы
  % развёртывания для нескольких наборов криптографических параметров.}

  \workproject{\textbf{Приватная агрегация для федеративного обучения с подкреплением}}
  \workprojectdetail{Участвовал в разработке TWAPPA-FRL, объединяющей
  доверительно-взвешенную агрегацию (TWA) и гомоморфную агрегацию с сохранением
  конфиденциальности (PPA) для совместного обнаружения вторжений,
  в качестве автора с равным вкладом.}
  % \workprojectdetail{Объединил взвешивание клиентов на основе доверия
  % с зашифрованной CKKS-агрегацией обновлений федеративной модели и расшифрованием
  % только агрегированного результата отдельным центром управления ключами,
  % не участвующим в обучении.}
  \workprojectdetail{Разработал модель угроз и анализ безопасности, включая
  конфиденциальность индивидуальных обновлений, явно определённую утечку,
  а также предположения и ограничения путей выполнения с PPA.}

  \workproject{\textbf{Верифицируемое многоключевое шифрование с возможностью поиска}}
  \workprojectdetail{Участвовал в разработке EVMK-SSE, эффективной схемы
  многоключевого шифрования с возможностью поиска на основе забывчивой передачи
  и ослабленной OPRF, не зависящей от клиента.}
  \workprojectdetail{Проверил анализ безопасности, охватывающий конфиденциальность
  запросов, явную утечку шаблонов поиска и доступа, а также верифицируемость
  результатов при наличии вредоносного облачного сервера.}
}

\vspace{0.50em}

\cvrow{}{%
  \workentry
    {\textbf{Стажёр ИБ (EDR, TI, SOC)} @
    \href{https://global.ptsecurity.com/en/}{Positive Technologies}}
    {авг. 2024--февр. 2025}

  \workproject{Разработал на Python инструменты с unit-покрытием для
  проверки путей Windows, сетевых конфигураций и политик контроля конечных устройств.}
  \workproject{Построил конвейер обработки вредоносного ПО на базе Airflow:
  сбор образцов из VX-Underground, хранение артефактов в S3/MinIO
  и сканирование правилами YARA.}
  \workproject{Проводил анализ EVTX, Powershell и вредоносные VBA-макросы;
  извлекал IoC и сопоставлял TTP с MITRE ATT\&CK для CVE-2020-5902, CVE-2023-38831.}
}

\vspace{0.50em}

\cvrow{}{%
  \workentry
    {\textbf{Инженер-программист} @ \href{https://academai.ru/}{AcademAI}}
    {авг. 2023--авг. 2024}

  \workproject{Реализовал REST API на Python и конвейер генерации на базе LangChain
  с использованием структурированных промптов и асинхронной оркестрации LLM.}
  \workproject{Разработал слой приложения на Next.js, TypeScript,
  Prisma и Auth.js, включая аутентификацию и постоянное хранение данных курсов.}
  \workproject{Настроил CI/CD с использованием Docker и GitHub Actions.}
}

\vspace{0.75em}

\cvrow{Образование}{%
  \raggedright

  % ── Магистратура ────────────────────────────────
  \textbf{Чунцинский университет почты и телекоммуникаций},
  Чунцин, Китай\par
  \textit{Школа кибербезопасности и информационного права}\par
  \vspace{0.12em}

  \eduentry
    {$\diamond$ \textbf{Магистратура} по направлению «Сетевая и информационная безопасность»}
    {сент. 2024--июль 2027 \textit{(ожидается)}}

  \edudetail{Научный руководитель:
    \href{https://faculty.cqupt.edu.cn/tangfei/en/index.htm}{проф. Тан Фэй}}
  \edudetail{Средний балл: 3,98/4,00}

  \vspace{0.48em}

  % ── Бакалавриат ─────────────────────────────────
  \textbf{Ухтинский государственный технический университет},
  Ухта, Россия\par
  \vspace{0.12em}

  \eduentry
    {$\diamond$ \textbf{Бакалавриат} по направлению «Информационные системы и технологии»}
    {март 2022--июнь 2024}

  \edudetail{Средний балл: 4,49/5,00}

  \vspace{0.48em}

  % % ── Предыдущее обучение в бакалавриате ─────────
  % \textbf{Санкт-Петербургский государственный университет телекоммуникаций
  % им. проф. М.\,А. Бонч-Бруевича},
  % Санкт-Петербург, Россия\par
  % \vspace{0.12em}

  % \eduentry
  %   {$\diamond$ \textbf{Обучение в бакалавриате} по направлению
  %   «Информационные системы и технологии»}
  %   {сент. 2020--март 2022}
}

\vspace{0.75em}

\cvrow{Публикации}{%
  \pub{\underbar{Artur Iasenovets}, Fei Tang, H. Zhu, P. Wang, and L. Liu.}
  {``\href{https://arxiv.org/abs/2511.02656}
  {Bringing Private Reads to Hyperledger Fabric via Private Information Retrieval},''
  arXiv:2511.02656, 2025. Рукопись находится на рецензировании.}

  \pub{H. Zhu, F. Tang, W. Qin, J. Shan, P. Wang, and \underbar{Artur Iasenovets}.}
  {``\href{https://doi.org/10.1007/s44443-025-00405-8}
  {EVMK-SSE: Efficient and Verifiable Multi-Keyword Symmetric Searchable Encryption},''
  \textit{Journal of King Saud University -- Computer and Information Sciences}, 2025.}

  \pub{\underbar{Artur Iasenovets} and Fei Tang.}
  {``\href{https://doi.org/10.60797/IRJ.2026.163.57}
  {MFCTIIF: Multi-Feed Cyber Threat Intelligence Integration Framework},''
  \textit{International Research Journal}, № 1 (163), 2026.}

  % \vspace{0.15em}
  % \textbf{Рукопись в подготовке}\par\vspace{0.25em}

  % \pub{Piyumi M. S. Rajapaksha$^{*}$, \underbar{Artur Iasenovets$^{*}$},
  % Yinxue Yi, and Fei Tang.}
  % {``TWAPPA-FRL: Trust-Weighted Privacy-Preserving Aggregation for Federated
  % Reinforcement Learning in Collaborative Intrusion Detection,''
  % рукопись в подготовке, 2026. $^{*}$Равный вклад.}
}

\vspace{0.75em}

\cvrow{Награды и\\достижения}{%
  \entry
    {\textbf{Стипендия правительства КНР (CSC) для обучения в магистратуре}}
    {2024--2027}
  Полная стипендия на весь период обучения в магистратуре CQUPT.

  \vspace{0.45em}
  \entry
    {\textbf{Победитель конкурса EdTech-проектов с проектом YouKnow}}
    {2023}
  Победа в конкурсе, организованном Физтех-школой МФТИ,
  Фондом развития Физтех-школ и Startech.Base.

  \vspace{0.45em}
  \entry
    {\textbf{Диплом I степени, «Севергеоэкотех»}}
    {2022}
  Первое место в секции «Современные информационные технологии»
  XXIII Международной молодёжной научной конференции.

  \vspace{0.45em}
  \entry
    {\textbf{Региональный демо-день Sber Student Accelerator}}
    {2023}
  Завершил программу Sber Student Accelerator в составе команды AcademAI
  и принял участие в региональном демо-дне Северо-Западного округа.

  \vspace{0.3em}
  \href{https://iasenovets.github.io/files/certificates_selected_en_public_ready_v2.pdf}
       {\small Подтверждающие сертификаты}
}

\end{document}
